\section{Implementation in the BX3 Framework}
\label{sec:implementation}

The Bailout Protocol is not an isolated mechanism — it is the \emph{runtime expression}
of the BX3 Framework's upstream accountability guarantee. Where the three-layer
architecture (Purpose Layer, Bounds Engine, Fact Layer) defines the static structure
of accountability, the Bailout Protocol operationalises that structure as a living,
deterministic state machine that fires whenever a node exhausts its local resolution
capacity. This section specifies how the Protocol integrates with each layer of BX3,
how the Bounds Engine triggers bail-out conditions, how the Fact Layer enforces the
immutable audit trail, and how the state machine itself functions as an extension of
the Fact Layer.

% ─────────────────────────────────────────────
% 4.1  Integration with the Purpose Layer
% ─────────────────────────────────────────────

\subsection{Integration with the Purpose Layer}
\label{sec:purpose-layer}

The Purpose Layer is the human accountability anchor of the BX3 Framework. It defines
intent, scope, constraints, and — critically — the identity and authority of every human
principal who bears upstream accountability for system outcomes. The Bailout Protocol
integrates with the Purpose Layer through two distinct mechanisms: the \emph{anchor
registry} and the \emph{directive binding}.

\bigskip

\noindent\textbf{Anchor Registry.}

Every human principal is registered in the Purpose Layer's \texttt{RoleDefinitionLanguage}
(RDL) schema with an \texttt{ActorType = HUMAN} classification and a
\texttt{ScopeOfAuthority} field that defines the boundary of their accountability domain.
When a node initiates an escalation, it queries the anchor registry via
\texttt{GetNextUpstreamAnchor}($n$) to determine the nearest upstream human-occupied
Purpose Layer node. The registry is maintained by the Purpose Layer and is immutable
at runtime — a human principal's authority scope cannot be modified by any Bounds Engine
or Fact Layer component; changes require a Purpose Layer update transaction signed by
the principal themselves.

\bigskip

\noindent\textbf{Directive Binding.}

When a human principal issues a resolution directive during the
\texttt{HUMAN\_ACKNOWLEDGED} state, that directive is bound to the trigger package and
propagated back down the escalation chain to the originating node. The directive carries
the principal's RDL identity, a cryptographic signature over the directive payload, and
a timestamp with non-repudiable provenance. The Fact Layer records the directive as the
definitive resolution of the trigger. No machine actor may modify, supersede, or withdraw
a human directive once issued; if the directive cannot be executed, the node transitions
to \texttt{RESOLVED\_EXECUTION\_FAILURE} and the Fact Layer logs the discrepancy for
post-hoc review.

\bigskip

\noindent\textbf{Human Accountability Anchor Properties.}

The Purpose Layer enforces three structural properties that make human anchors
capable of terminating escalation chains:

\begin{enumerate}
\item \textbf{Legal authority:} A human principal's directives carry legal weight that
  machine actors cannot replicate. Only a human can authorise financial commitments,
  regulatory determinations, or contractual obligations.

\item \textbf{Scope fidelity:} The \texttt{ScopeOfAuthority} field ensures a human is
  never asked to resolve exceptions outside their domain competence. An escalation that
  reaches a human anchor has by definition reached an authority class capable of binding
  the organisation.

\item \textbf{Non-repudiability:} Human directives are cryptographically signed and
  timestamped at the Purpose Layer. A human cannot later disclaim knowledge of a
  directive they issued — the signature is embedded in the trigger package's immutable
  record.
\end{enumerate}

% ─────────────────────────────────────────────
% 4.2  Bounds Engine Triggering Bail-Out
% ─────────────────────────────────────────────

\subsection{Bounds Engine Triggering Bail-Out}
\label{sec:bounds-trigger}

The Bounds Engine is the reasoning component of the BX3 Framework. It maintains the
\texttt{BoundsTable}, evaluates proposed actions against all registered constraints,
and determines whether a given action can proceed or must be escalated. The Bailout
Protocol is triggered exclusively through the Bounds Engine's \texttt{BailoutCondition}
evaluation, which runs continuously during normal operation.

\bigskip

\noindent\textbf{BailoutCondition Evaluation.}

The Bounds Engine evaluates $\texttt{BailoutCondition} \equiv \bigvee_{i=1}^{5} C_i$
continuously at each sensing cycle. Each condition $C_i$ maps to a distinct failure
mode:

\bigskip

{\setlength{\tabcolsep}{6pt}
\begin{tabular}{lp{0.62\linewidth}}
\textbf{$C_1$ — ConstraintViolation} & A proposed action violates one or more
constraints in the \texttt{BoundsTable} and the violation cannot be resolved by
constraint relaxation within the node's authority class. The Bounds Engine logs the
specific constraint that failed and the proposed relaxation attempt. \\
\textbf{$C_2$ — NovelInput} & The input state $x$ falls outside the training
distribution $D_{\text{train}}$ by a margin exceeding the node's novelty threshold
$\tau_{\text{novelty}}$, and no verified response strategy exists in the
\texttt{ResponseTable}. The Bounds Engine records the Mahalanobis distance from the
distribution centre and the ResponseTable miss. \\
\textbf{$C_3$ — ResourceExhaustion} & A required resource (time, memory, API quota,
external dependency) is depleted such that safe completion of the current task cannot
be guaranteed. The Bounds Engine logs resource utilisation percentages at time of
trigger. \\
\textbf{$C_4$ — HumanDirectiveConflict} & Two or more directives from distinct
Purpose Layer anchors express conflicting intent, and the node lacks the authority to
adjudicate. The Bounds Engine logs both directives and the specific axis of conflict. \\
\textbf{$C_5$ — ForensicLedgerFailure} & A write to the Forensic Ledger fails or does
not return a confirmation within the designated confirmation window. This condition is
checked at every write operation; failure triggers immediate escalation because the
audit trail cannot be guaranteed.
\end{tabular}}

\bigskip

When any $C_i$ evaluates to \texttt{True}, the Bounds Engine transitions the state
machine to \texttt{BOUNDED} and begins the local resolution procedure. The trigger
package $\mathcal{T}$ assembled at this point contains the condition sub-type, the node
identifier, the full state snapshot $\mathbf{s}_9$, and the current confidence score
$\phi$. The package is immutable from the moment of assembly — no machine actor may
modify the trigger package after it enters the \texttt{BAIL\_PENDING} state.

\bigskip

\noindent\textbf{Local Resolution and Transition to \texttt{BAIL\_PENDING}.}

During \texttt{BOUNDED}, the Bounds Engine searches the action space $A$ for any action
$a$ satisfying $\texttt{isPermissible}(a) \land \texttt{isSafe}(a) \land \texttt{hasResource}(a)$.
If such an $a$ exists, the node resolves the exception locally and transitions to
\texttt{IDLE}. If no such $a$ exists — formally, $\neg \exists \; a \in A :
\texttt{isPermissible}(a) \land \texttt{isSafe}(a) \land \texttt{hasResource}(a)$ — the
Bounds Engine transitions to \texttt{BAIL\_PENDING}, assembling the trigger package for
upstream propagation. The local resolution procedure is time-bounded by $T(\texttt{BOUNDED})
= 30$ seconds; if the timeout fires before resolution, the state transitions
automatically to \texttt{BAIL\_PENDING} regardless of whether a resolution candidate
has been found.

% ─────────────────────────────────────────────
% 4.3  Fact Layer Enforcing Audit Trail
% ─────────────────────────────────────────────

\subsection{Fact Layer Enforcing Audit Trail}
\label{sec:fact-layer}

The Fact Layer is the enforcement and audit component of the BX3 Framework. It is
responsible for deterministic action execution, cryptographic sealing of all state
transitions, and the immutable maintenance of the Forensic Ledger. The Bailout
Protocol is deeply integrated with the Fact Layer in three ways: state machine
enforcement, cryptographic integrity, and tamper-evident chain construction.

\bigskip

\noindent\textbf{State Machine as Fact Layer Extension.}

The Bailout Protocol state machine $\mathcal{S}_{\text{bailout}}$ is not a separate
subsystem — it is an extension of the Fact Layer's enforcement authority. The Fact
Layer owns the state transition function $\delta$; no Bounds Engine or Purpose Layer
component may invoke a state transition directly. When the Bounds Engine evaluates a
trigger condition and determines that escalation is required, it \emph{requests} a state
transition from the Fact Layer. The Fact Layer validates the transition request against
the current state and event, executes the transition atomically, and seals the result
in the Forensic Ledger before acknowledging the transition to the requesting component.

This enforcement pattern ensures that the state machine cannot be manipulated by
any component outside the Fact Layer. A machine actor that attempts to skip states,
forge a transition, or fabricate a trigger condition will fail at the Fact Layer's
validation step, producing a $C_5$ (\texttt{ForensicLedgerFailure}) condition that
triggers its own escalation.

\bigskip

\noindent\textbf{Cryptographic Sealing.}

Every state transition in $\mathcal{S}_{\text{bailout}}$ produces a Forensic Ledger
entry with the following structure:

\begin{equation}
\mathcal{L}_{\text{bailout}} = \bigl\langle
  \tau_{\text{id}},           % Trigger ID (globally unique)
  q_{\text{from}},            % Source state
  q_{\text{to}},              % Target state
  e,                          % Triggering event
  n_{\text{actor}},           % Node executing transition
  t_{\text{transition}},      % Wall-clock timestamp (nanosecond precision)
  \mathbf{s}_9,              % State snapshot at transition
  \phi,                       % Confidence score at fire time
  \text{digest},             % SHA-256 digest: H(this\_entry \Vert previous\_digest)
  \rho_{\text{final}}         % Final resolution (populated at RESOLVED)
\bigr\rangle
\end{equation}

The digest field implements a chain integrity mechanism. Each entry's digest is computed
as the SHA-256 hash of the concatenation of the entry's raw fields and the previous
entry's digest. This creates a tamper-evident chain: modifying any historical entry
breaks the digest chain from that point forward, and the break is detectable by any
independent auditor without access to the chain's internal state.

\bigskip

\noindent\textbf{Audit Completeness Enforcement.}

The Fact Layer enforces the Audit Completeness Guarantee through a mandatory write
confirmation requirement. A state transition is not considered complete until the
Forensic Ledger has confirmed the write with a positive acknowledgment. If the write
fails or the acknowledgment is not received within the confirmation window, the Fact
Layer sets $C_5$ (\texttt{ForensicLedgerFailure}) and triggers an autonomous escalation
from the current state. The node does not proceed with any action that depends on an
unconfirmed state transition.

% ─────────────────────────────────────────────
% 4.4  Forensic Ledger Recording
% ─────────────────────────────────────────────

\subsection{Forensic Ledger Recording}
\label{sec:forensic-ledger}

Every Bailout Protocol event produces one or more Forensic Ledger entries across three
planes, corresponding to the BX3 Framework's three-layer architecture:

\bigskip

\begin{trivlist}\itemindent2em
\item[\textbf{Purpose Plane}] Records authorization to escalate, human acknowledgment
  events, directive issuance, and principal identity at each escalation hop. These
  entries establish who knew about the exception and what they directed.

\item[\textbf{Bounds Engine Plane}] Records trigger condition evaluations ($C_1$ through
  $C_5$), local resolution attempts, constraint violation details, confidence scores
  at fire time, and escalation path metadata. These entries establish what the machine
  observed and why it concluded escalation was necessary.

\item[\textbf{Fact Plane}] Records state transition timestamps, timeout fire events,
  transition validation confirmations, resolution action executions, and Forensic
  Ledger write acknowledgments. These entries establish what the system did at each
  step and that the audit trail itself is intact.
\end{trivlist}

\bigskip

The Forensic Ledger entries for a complete trigger lifecycle form a directed acyclic
graph rooted at the originating node's first \texttt{BOUNDED} entry. The graph
captures the full provenance of the exception: initial condition, local resolution
attempt, escalation path (all hops), human acknowledgment(s), directive (if issued),
actuator execution (or execution failure), and closure. No entry in this graph can
be deleted, modified, or reordered after commitment. The chain is append-only and
cryptographically sealed.

\bigskip

\begin{note}
\textbf{Retention and Access Control.} Forensic Ledger entries for Bailout Protocol
events are retained for the full operational lifetime of the system. Access to read
the ledger is role-scoped via the Purpose Layer: auditors may read all entries,
operators may read entries for nodes in their domain, and no machine actor may read
entries for nodes outside its own trust boundary without an explicit Purpose Layer
authorization grant.
\end{note}

% ─────────────────────────────────────────────
% 4.5  Bailout Protocol as Runtime Expression of Upstream Accountability
% ─────────────────────────────────────────────

\subsection{Bailout Protocol as Runtime Expression of Upstream Accountability}
\label{sec:runtime-expression}

The BX3 Framework's upstream accountability guarantee is a static architectural
commitment: every node has a guaranteed path to a human principal who bears
accountability for the node's outcomes. The Bailout Protocol is the mechanism that
makes this commitment operational — it is the runtime expression of the upstream
accountability guarantee, the bridge between the static promise of human oversight
and the dynamic reality of autonomous operation.

\bigskip

The relationship can be stated formally as a theorem:

\bigskip

\noindent\textbf{Theorem (Upstream Accountability Operationalisation).}
\textit{For any node $n$ in a BX3 system, if the Bounds Engine evaluates
$\texttt{BailoutCondition} \equiv \text{True}$ and local resolution fails, then the
Bailout Protocol guarantees that the trigger package reaches a Human Accountability
Anchor within a finite number of hops, or the system halts with a timeout exception
and enters a \texttt{HALT\_REQUIRED} state awaiting human review.}

\bigskip

\noindent\textit{Proof sketch.} The trigger condition $C_i$ ensures that any
condition a Bounds Engine cannot resolve locally is captured and classified. The
bypass rule $\forall n \in \text{Nodes} : \texttt{IsMachineActor}(n) \implies
\texttt{CannotBeEscalationEndpoint}(n)$ guarantees that the escalation traversal never
terminates at a machine actor. The finite state machine has a single terminal state
(\texttt{RESOLVED}) and no dead-end states. If a human is reachable within
\texttt{MaxEscalationDepth} hops, the escalation terminates there. If not, the timeout
at \texttt{MaxEscalationDepth} drives the state to \texttt{RESOLVED} as a timeout
exception with \texttt{HALT\_REQUIRED} set, ensuring the node cannot resume autonomous
operation. $\square$

\bigskip

This theorem establishes that the Bailout Protocol is not a best-effort safety
feature — it is a deterministic guarantee enforced by the structural properties of
the state machine, the bypass rule, and the Fact Layer's enforcement authority. The
guarantee holds regardless of network conditions, node failures, or the number of
machine actors in the escalation path.

\bigskip

The upstream accountability guarantee also operates in reverse: when a human principal
issues a directive, that directive propagates back down the same path with the same
immutability guarantees, ensuring that the principal's intent is executed faithfully
at the originating node. This bidirectional accountability chain — escalation on
exception, resolution on directive — is the operationalisation of the BX3 Framework's
foundational commitment to never orphan an action from human responsibility.

% ─────────────────────────────────────────────
% 4.6  Formal Properties Summary
% ─────────────────────────────────────────────

\subsection{Summary of Formal Properties}
\label{sec:properties-summary}

The implementation of the Bailout Protocol within the BX3 Framework satisfies the
following formal properties:

\bigskip

\begin{trivlist}\itemindent0em
\item[\textbf{Upstream Accountability Guarantee}] For any node $n$, there exists a
  finite escalation path $P = \langle n = p_0, p_1, \ldots, p_k \rangle$ such that
  $p_k$ is a Human Accountability Anchor, and $\forall i \in [0, k-1]$,
  $\texttt{IsMachineActor}(p_i) = \text{True}$. The path is constructed by the
  \texttt{GetNextUpstreamAnchor} algorithm operating on the Purpose Layer anchor
  registry.

\item[\textbf{Non-Orphaning Guarantee}] No trigger can reach \texttt{RESOLVED} without
  at least one human acknowledgment in its escalation path, unless the resolution is a
  timeout exception at \texttt{MaxEscalationDepth}, in which case the node enters
  \texttt{HALT\_REQUIRED} state and cannot resume autonomous operation without human
  clearance.

\item[\textbf{Deterministic State Transitions}] Every transition in
  $\mathcal{S}_{\text{bailout}}$ is deterministic and enforced by the Fact Layer.
  No machine actor may invoke, skip, or modify a state transition.

\item[\textbf{Audit Completeness Guarantee}] Every state transition produces a
  confirmed Forensic Ledger entry across all three planes (Purpose, Bounds Engine,
  Fact) before the transition is considered complete. A transition without a
  confirmed write is invalid and triggers $C_5$.

\item[\textbf{Immutable Trigger Package}] After the Bounds Engine assembles the
  trigger package $\mathcal{T}$ at \texttt{BOUNDED}, no machine actor may modify any
  field of $\mathcal{T}$. Modifications are detectable at Fact Layer validation and
  produce a forensic anomaly entry.

\item[\textbf{Graceful Degradation}] Network disconnection does not suppress or delay
  escalation. Triggers generated during disconnection are buffered at the originating
  node and delivered upon reconnection, preserving the full escalation path in the
  Forensic Ledger.
\end{trivlist}

\bigskip

These properties collectively ensure that the Bailout Protocol, as implemented within
the BX3 Framework, delivers the BX3 Framework's core safety promise in all operational
contexts: \emph{accountability always escalates to a human, or the system halts and
awaits human review.}